\settrack{trackbridge}
% VOICE: 3rd-person analytical/reconstructive

\chapter{Growing a Mind}
\label{pos14:growing-a-mind}

% Source: manuscript/sources/turing-death.md (imported Plan 0011)

\begin{quote}
\textit{``A computer would deserve to be called intelligent if it could deceive a human into believing that it was human.''}

\hfill --- Alan Turing
\end{quote}
\vspace{0.5cm}

\srcnote{turing-death.md | manuscript/bridge/pos14-growing-a-mind.tex | Turing morphogenesis thread, Kauffman connection}

The question of how biological systems compute --- not what they compute, but the physical substrate and developmental process by which computational structure arises --- occupied Alan Turing in the final years of his life. It is a question that would eventually lead, by a circuitous route, to the central proposition of this book.

\section*{The Mathematician}

Alan Turing laid the theoretical foundations of digital computation in 1936.\footcite{turing1936} His paper that year established the abstract model now called a Turing machine --- a formalism that defines what is and is not computable, independent of any physical implementation. During the Second World War, Turing applied these ideas at Government Communications Headquarters, where his contributions to breaking German military ciphers are estimated by historians to have shortened the European war by more than two years and saved millions of lives.\footcite{winterbotham1974} The work remained classified for decades.

By the time transistors began displacing vacuum tubes in the early 1950s --- Bardeen and Brattain's point-contact device at Bell Labs in 1947 had opened that path\footcite{bardeen1948} --- Turing had turned his attention elsewhere. The question that absorbed him was not how to make faster digital machines, but whether digital machines represented the only viable path to artificial intelligence at all.

\section*{Morphogenesis as Computation}
%TODO: ILLUSTRATION CANDIDATE. This section could benefit from an image showing a cellular automaton pattern alongside an identical pattern on a natural sea shell (several examples in Wolfram's NKS). The book is intentionally all-text, but this is the one place where a picture genuinely teaches — it shows the reader that computation-as-physics is not metaphor, it's literal. The CA and the shell run the same algorithm. Discuss with Bruce: source/licensing (NKS images), placement, whether it breaks the text-only aesthetic or earns its place.

Turing's final research program addressed morphogenesis: the biological process by which organisms develop their characteristic forms. The program was motivated by a specific observation. A brain is not designed and assembled component by component. It grows. The computational structure of a brain is an emergent property of a developmental process, shaped by genetics and environment in ways that no engineer could replicate by direct fabrication. If artificial intelligence was the target, Turing reasoned, the developmental pathway deserved as much attention as the end state.

D'Arcy Wentworth Thompson had published the foundational work in mathematical biology decades earlier,\footcite{thompson1917} but the field remained underdeveloped. Turing's contribution --- the reaction-diffusion model published in 1952 as ``The Chemical Basis of Morphogenesis''\footcite{turing1952} --- provided a mathematical framework for how chemical gradients could spontaneously generate spatial patterns in biological tissue. The model demonstrated that uniform conditions could, under the right parameters, self-organize into structured differentiation. A brain, on this account, is not assembled. It precipitates.

The proposition is not that Turing solved the morphogenesis problem. He did not. The proposition is that he asked the right question and established the mathematical vocabulary in which subsequent researchers would pursue it.

\section*{The End of the Work}

% NOTE: The apple was never tested for cyanide. The inquest ruled suicide by cyanide
% poisoning, but the delivery mechanism is debated. See source-facts.md and Plan 0008 Fix m1.

In 1952, Turing was convicted under British law for gross indecency following his admission of a consensual homosexual relationship. The court offered a choice between imprisonment and chemical castration; he accepted the latter. The treatment entailed regular injections of synthetic estrogen, with significant physiological side effects, and continued for the duration of his sentence. On 8 June 1954, Turing was found dead at his home in Wilmslow. The inquest returned a verdict of suicide by cyanide poisoning. The precise circumstances remain disputed; the outcome does not.

Turing's wartime contributions were not publicly known at the time of his death. They would not be declassified until the mid-1970s.

\section*{The Thread Continues}
%TODO: REWRITE — Bruce found this section too dense to read. Four consecutive paragraphs of academic exposition with no concrete imagery. Specific problems:
% 1. Kauffman "catalytic closure" introduced without plain-language anchor
% 2. McCulloch-Pitts citation (line 50) drops cold — adds complexity, no GA payoff. Cut or fold into a sentence.
% 3. "brittle in a specific sense" / "regime not reachable by engineering to specification" — academic register needs concrete analogies
% 4. Three frameworks stacked (Turing → Kauffman → Wolfram) with no breathing room
% 5. Summary paragraph (line 56) restates everything abstractly a second time
% Approach: lead each idea with something concrete (bathtub, ant colony, growing a garden vs building a house), then state the abstraction. Reduce citation density. Give the reader one idea per paragraph, not three. Keep the content — the logic chain is correct — but soften the delivery.

The morphogenesis program did not end with Turing. From the 1960s onward, a cohort of researchers extended his framework --- asking not just how organisms develop their shape, but how sufficiently complex autocatalytic chemistry might give rise to self-sustaining computational structure. Stuart Kauffman's work on autocatalytic sets, developed through the 1970s and 1980s,\footcite{kauffman1993} supplied the missing link. Kauffman demonstrated mathematically that above a threshold of molecular diversity, catalytic closure becomes overwhelmingly probable. Self-sustaining chemical networks, networks that maintain and reproduce their own structure, can emerge spontaneously from sufficiently complex mixtures.

The implication for computation is direct. If Turing was correct that biological intelligence arises through developmental processes rather than assembly, and if Kauffman is correct that developmental processes are a special case of autocatalytic emergence, then the path to artificial intelligence may run through chemistry and self-organization rather than through engineering and fabrication. According to this account, a mind is not built. It is grown.

McCulloch and Pitts had formalized neural network mathematics as early as 1943,\footcite{mccullochpitts1943} but the question of how a physical system acquires the topology of a neural network remained open. Kauffman's framework suggests an answer: the topology emerges at criticality, at the edge of chaos, as an autocatalytic system self-organizes under the right conditions. The network does not have to be designed. It has to be tuned.

The distinction matters practically. An engineered neural network is brittle in a specific sense: its capabilities are bounded by the topology its designers specify. A grown neural network, one whose topology is a product of developmental dynamics, can in principle acquire structure that no designer anticipated. The distinction between training a static architecture and allowing an architecture to evolve during training is not merely technical; it is a difference in kind. Kauffman's edge-of-chaos result suggests that systems operating near criticality maximize both order and adaptability simultaneously --- a regime not reachable by engineering to specification.

Stephen Wolfram's subsequent work on cellular automata and computational universality\footcite{wolfram2002} reinforced the picture: simple rule-governed systems can generate behavior of arbitrary complexity. The morphogenesis question, reformulated in computational terms, becomes: what minimal conditions are required for a physical system to spontaneously generate computational universality? Turing's reaction-diffusion equations, Kauffman's autocatalytic sets, and Wolfram's universality results are three successive answers to variants of the same question.

This is the conceptual bridge that connects Turing's morphogenesis program to the central proposition of the book. The proposition holds that a quantum substrate --- a two-dimensional electron gas operating at the edge of quantum criticality --- can undergo a process formally analogous to Kauffman's autocatalytic emergence, producing neural network topology as a spontaneous consequence of dynamics rather than as a product of engineering. Whether that process has actually been realized is the question the subsequent chapters address.

\chapterreturn
